\chapter{The Pattern}

\vspace{1em}
\hfill
\begin{minipage}[c]{0.6\textwidth}
\raggedleft
\small\itshape
[The self is] nothing but a bundle or collection of different perceptions, which succeed each other with an inconceivable rapidity, and are in a perpetual flux and movement.\\[0.5em]
\upshape\footnotesize David Hume, \textit{A Treatise of Human Nature} (1739)
\end{minipage}%
\hspace{0.5em}
\begin{minipage}[c]{0.12\textwidth}
% \includegraphics[width=\textwidth]{figures/ch14_pattern.png}
\end{minipage}
\vspace{1.5em}

\noindent \textsc{The word \emph{you} has done a great deal of work} in this book, and it is time to look at what it names.

\section*{The temporal dissolution, recalled}

\textit{Worldlines} already took the self apart once, along the time axis. A person is a worldline, a four-dimensional shape, and the five-year-old and the eighty-year-old are equally real cross-sections of it. There is no thin moving spark of selfhood travelling down the curve; there is the whole curve, and \emph{I} is the word a cross-section uses for the shape it belongs to. The self was not destroyed by this. It was relocated, from a moving point to a static pattern.

This book takes the self apart a second time, along a different axis.

\section*{The self as a local model}

Recall the working model from Chapter 8: an observer builds tractable local models of its surroundings, including a model of itself. We can represent that self-model by a program $p_{\text{self}}$, a compressed and useful approximation of the person and situation. It helps the system predict its next move, hold a history, and say \emph{I}. It is not a further thing seated inside the observer. It is a description, not the whole of what it describes.

Two thinkers reached this without the Library. Derek Parfit argued that personal identity is not a deep further fact: what there is, is a stream of mental states with more or less connectedness, and ``the same person'' is a way of summarizing the connectedness, not a thread running underneath it. Thomas Metzinger argued that the self is a model the brain runs and cannot see through, a representation so transparent the system mistakes it for a thing. Both land where Chapter 8 landed: the self is a compression, useful, finite, and not the structure it compresses.

\section*{The dissolution across descriptions}

Now place the model beside the Library.

Many program descriptions can encode patterns matching parts of your self-model. $M$ weights the programs, not observer occurrences. If Picture A supplies multiple matching occurrences, the label \emph{I} does not by itself choose among them, and the book supplies no self-location distribution. In ordinary life, the label still picks out the locally continuing person who remembers, acts, and bears responsibility.

And here the deepest thread of the book closes. The questions \emph{who am I}, \emph{where am I}, and \emph{which system am I modeling} all require boundaries. The self is a compression an observer uses; so is an index; so is a system treated as one definite thing. The postulate says none of those decompositions is view-independently canonical. This does not prove that many selves or worlds exist. It shows the same descriptive act at three levels. This is the convergence Chapter 2 promised: the cut met three times.

\section*{What is preserved}

This is not annihilation, and it is worth being exact about why.

Parfit's move was to separate identity from what matters. You need not posit a further metaphysical entity beyond psychological and bodily continuity to preserve memories, commitments, responsibility, and love. Identity may not have been doing all the work assigned to it. The pattern and its connections do that work here.

Carry that into the Library and one conclusion holds without an occurrence census: no thread supplied by $M$ ties ``you here'' to a hypothetical ``you there.'' What matters locally, the features, connectedness, and braid, does not require such a thread. If this lands as vertigo rather than relief, hold it lightly: that the self is a cut is a claim about a boundary and a model, not a claim that the person or experience is unreal. You are here. What dissolves is the demand that the line be view-independently canonical.

\section*{Why this lets care be care}

There is a consequence the last chapter was waiting for.

If the self is a convention, a label drawn around a region, then so is the person you love, and so is the braid between you: a relation between two conventions. This sounds as though it should be the final dilution, the point where care evaporates into bookkeeping about labels. It is the reverse.

A relationship is not maintained by metaphysical labels alone. Its connectedness persists through memory, attention, repair, and repeated acts. The framework has no stake in it. Care is therefore a practice, performed here, rather than a guarantee the cosmos underwrites. The dissolution of a canonical self helps turn care from a thing you possess into a thing you do.

That is the last piece. The book can now say what it has been building toward.

The next chapter is the last.
